\chapter{Lesson 4: Family}

\section{Vocabulary}

\begin{center}
\begin{tabular}{llll}
\toprule
Characters & Pinyin & Part of speech & Meaning \\
\midrule
\zh{家} & jiā & noun & family, home \\
\zh{爸爸} & bàba & noun & dad \\
\zh{妈妈} & māma & noun & mom \\
\zh{哥哥} & gēge & noun & older brother \\
\zh{姐姐} & jiějie & noun & older sister \\
\zh{弟弟} & dìdi & noun & younger brother \\
\zh{妹妹} & mèimei & noun & younger sister \\
\zh{儿子} & érzi & noun & son \\
\zh{女儿} & nǚ'ér & noun & daughter \\
\zh{孩子} & háizi & noun & child \\
\zh{有} & yǒu & verb & to have; there is \\
\zh{没} & méi & adverb & not (only with \zh{有} and past actions) \\
\zh{几} & jǐ & question word & how many (small number) \\
\zh{个} & gè & measure word & general measure word \\
\zh{口} & kǒu & measure word & for members of a family \\
\zh{两} & liǎng & numeral & two (of something) \\
\zh{和} & hé & conjunction & and (between nouns) \\
\zh{都} & dōu & adverb & all, both \\
\zh{这} / \zh{那} & zhè / nà & pronoun & this / that \\
\zh{谁} & shéi & question word & who \\
\zh{照片} & zhàopiàn & noun & photograph \\
\zh{张} & zhāng & measure word & for flat things \\
\zh{的} & de & particle & possessive marker \\
\zh{医生} & yīshēng & noun & doctor \\
\zh{律师} & lǜshī & noun & lawyer \\
\zh{工程师} & gōngchéngshī & noun & engineer \\
\zh{大学生} & dàxuéshēng & noun & university student \\
\zh{做} & zuò & verb & to do \\
\zh{工作} & gōngzuò & noun / verb & job; to work \\
\bottomrule
\end{tabular}
\end{center}

\section{The family tree}

Chinese kinship terms are much more specific than English ones: there is no single word
for ``brother'' or ``sister''; relative age is built in. This is my (made-up) family for
the oral presentation.

\begin{center}
\begin{tikzpicture}[p/.style={draw, rounded corners, align=center, font=\small, minimum width=1.9cm, minimum height=1cm}, m/.style={p, fill=blue!8}, f/.style={p, fill=red!8}]
  \node[f] (nn) at (-3.2,4) {\zh{奶奶}\\ nǎinai};
  \node[m] (yy) at (-1.0,4) {\zh{爷爷}\\ yéye};
  \node[m] (wg) at (2.2,4) {\zh{外公}\\ wàigōng};
  \node[f] (wp) at (4.4,4) {\zh{外婆}\\ wàipó};
  \draw (nn) -- (yy); \draw (wg) -- (wp);
  \node[m] (ba) at (-2.1,2) {\zh{爸爸}\\ bàba};
  \node[f] (ma) at (3.3,2) {\zh{妈妈}\\ māma};
  \draw (-2.1,4) -- (ba);
  \draw (3.3,4) -- (ma);
  \draw (ba) -- (ma);
  \draw (0.6,2) -- (0.6,1);
  \draw (-4.2,1) -- (5.4,1);
  \node[m] (gg) at (-4.2,0) {\zh{哥哥}\\ gēge};
  \node[f] (jj) at (-1.8,0) {\zh{姐姐}\\ jiějie};
  \node[p, very thick, fill=yellow!25] (me) at (0.6,0) {\zh{我}\\ wǒ};
  \node[m] (dd) at (3.0,0) {\zh{弟弟}\\ dìdi};
  \node[f] (mm) at (5.4,0) {\zh{妹妹}\\ mèimei};
  \foreach \n in {gg,jj,me,dd,mm} \draw (\n.north) -- ++(0,0.5);
  \node[font=\scriptsize, gray] at (-4.2,-0.85) {older};
  \node[font=\scriptsize, gray] at (-1.8,-0.85) {older};
  \node[font=\scriptsize, gray] at (3.0,-0.85) {younger};
  \node[font=\scriptsize, gray] at (5.4,-0.85) {younger};
  \node[font=\scriptsize, align=center] at (-2.1,4.95) {father's side};
  \node[font=\scriptsize, align=center] at (3.3,4.95) {mother's side (\zh{外} ``outside'')};
\end{tikzpicture}
\end{center}

\begin{aside}
Grandparents differ by side: \zh{爷爷}/\zh{奶奶} for the father's parents, and
\zh{外公}/\zh{外婆} (southern usage) or \zh{姥爷}/\zh{姥姥} (northern usage) for the
mother's. The \zh{外} ``outside'' reflects the old patrilineal view that a married
daughter joined her husband's family. Children also call any adult friend of the family
\zh{叔叔} \emph{shūshu} ``uncle'' or \zh{阿姨} \emph{āyí} ``auntie''.
\end{aside}

\section{Grammar}

\subsection{The possessive \zh{的}}

Noun/pronoun + \zh{的} + noun expresses possession: \zh{我的书} ``my book'',
\zh{老师的照片} ``the teacher's photo''.

\begin{important}[When \zh{的} is dropped]
With a personal pronoun before a family member, a close relationship or an institution
one belongs to, \zh{的} is normally omitted:
\zh{我爸爸}, \zh{你妹妹}, \zh{我家}, \zh{我们学校}. Using \zh{的} here is not wrong, just
less natural.
\end{important}

\subsection{\zh{有} and \zh{没有}}

\zh{有} means ``to have'' (and later ``there is''). Its negative is \zh{没有}, never
*\zh{不有}.
\begin{itemize}
  \item \zh{我有两个姐姐。} \emph{Wǒ yǒu liǎng ge jiějie.} I have two older sisters.
  \item \zh{他没有弟弟。} \emph{Tā méiyǒu dìdi.} He doesn't have a younger brother.
  \item \zh{你有哥哥吗？} / \zh{你有没有哥哥？} Do you have an older brother?
\end{itemize}

\subsection{Numbers and measure words}

In Chinese a number cannot directly modify a noun. A \emph{measure word} (\zh{量词}
liàngcí) goes in between:
\begin{center}
\textbf{number} (or \zh{这} / \zh{那} / \zh{几}) \ + \ \textbf{measure word} \ + \ \textbf{noun}
\end{center}

\begin{center}
\begin{tabular}{llll}
\toprule
Measure word & Used for & Example & Meaning \\
\midrule
\zh{个} gè & people, general & \zh{三个学生} & three students \\
\zh{口} kǒu & members of a family & \zh{我家有五口人} & my family has five people \\
\zh{张} zhāng & flat things & \zh{一张照片} & a photo \\
\zh{本} běn & books & \zh{两本书} & two books \\
\zh{位} wèi & people (polite) & \zh{这位老师} & this teacher \\
\zh{只} zhī & many animals & \zh{一只狗} & a dog \\
\bottomrule
\end{tabular}
\end{center}

\begin{important}[\zh{二} vs \zh{两}]
Before a measure word, ``two'' is \zh{两} liǎng: \zh{两个人}, not *\zh{二个人}. When
counting, in numbers like 12 or 20, and in ordinals, use \zh{二}: \zh{十二}, \zh{二十},
\zh{第二}.
\end{important}

When the noun is clear from context it can be dropped: \zh{这张照片是我的，那张是我哥哥的。}

\subsection{Questions about number: \zh{几}}

\zh{几} asks ``how many'' when the expected answer is small (under about ten). It needs a
measure word.
\begin{itemize}
  \item \zh{你家有几口人？}---\zh{我家有四口人：爸爸、妈妈、弟弟和我。}
  \item \zh{你有几个姐姐？}---\zh{我有一个姐姐。}
\end{itemize}
(For larger numbers we will use \zh{多少} \emph{duōshao} in Lesson 5.)

\begin{aside}
The enumeration comma \zh{、} (\zh{顿号} dùnhào) separates items in a list; \zh{和} joins
the last two. \zh{和} only joins nouns or pronouns: ``I study and work'' is \emph{not}
*\zh{我学习和工作}; use two clauses or \zh{也}.
\end{aside}

\subsection{The adverb \zh{都}}

\zh{都} ``all, both'' refers back to a plural subject and goes before the verb.
\begin{itemize}
  \item \zh{我爸爸妈妈都是医生。} My parents are both doctors.
  \item \zh{我们都不是律师。} None of us is a lawyer.
  \item \zh{我们不都是律师。} Not all of us are lawyers.
\end{itemize}

\begin{important}[\zh{都不} vs \zh{不都}]
\zh{都不} = ``none'' (all are not); \zh{不都} = ``not all''. With both \zh{也} and
\zh{都}, the order is \zh{也都}: \zh{他们也都是学生}.
\end{important}

\subsection{Demonstratives and \zh{谁}}

\begin{itemize}
  \item \zh{这是谁？}---\zh{这是我姐姐。} Who is this?---This is my older sister.
  \item \zh{这个男孩子是谁？} Who is this boy? (\zh{这} + MW + noun)
  \item \zh{这是谁的照片？} Whose photo is this?
\end{itemize}
In speech \zh{这} is often \emph{zhèi} and \zh{那} \emph{nèi} before a measure word.

\subsection{Asking about jobs}

\zh{你爸爸做什么工作？}---\zh{他是工程师。} ``What does your dad do?---He's an engineer.''
A common shorter form is \zh{你爸爸是做什么的？}

\section{Model description}

\begin{example}[Presenting a family photo]
\zh{这是我家的照片。我家有五口人：爸爸、妈妈、一个哥哥、一个妹妹和我。我爸爸是医生，我妈妈是律师。我哥哥是大学生，在滑铁卢大学学电脑。我妹妹很小，她九岁。我们都很喜欢狗，可是我们家没有狗。}

``This is a photo of my family. There are five of us \ldots\ My older brother is a
university student studying computers at Waterloo. My little sister is small; she's nine.
We all like dogs, but we don't have one.''
(\zh{滑铁卢} Huátiělú = Waterloo; \zh{可是} kěshì = but.)
\end{example}

\section*{Review}
\begin{itemize}
  \item Kin terms encode older/younger and father's/mother's side.
  \item \zh{的} marks possession; drop it with pronoun + close relation (\zh{我妈妈}).
  \item \zh{有} is negated only by \zh{没}: \zh{没有}.
  \item Number + measure word + noun; \zh{个} is general, \zh{口} counts family members.
  \item \zh{两} before measure words, \zh{二} in counting.
  \item \zh{几} + MW for small numbers; \zh{都} before the verb; \zh{都不} vs \zh{不都}.
\end{itemize}
